\begin{frame}[t]{Project Objectives}{What will the simulation answer?}
  \textbf{Obligatory:}
  \begin{itemize}
    \item \textbf{Carbon Intensity \& Temporal Trade-offs}\\
      Quantify how CO$_2$ savings depend on the pause threshold
      (gCO$_2$eq/kWh) versus added training time
      $\Rightarrow$ Pareto frontier for carbon-aware scheduling.
    \item \textbf{Impact Assessment on SOTA Architectures}\\
      Benchmark absolute and relative CO$_2$ reduction potential for
      different large models (e.g.\ DeepSeek V3, Kimi K2).
  \end{itemize}

  \vspace{0.3cm}
  \textbf{Wishful:}
  \begin{itemize}
    \item \textbf{Spatiotemporal Selection}\\
      Identify favourable \emph{time windows} and \emph{seasons}
      to run training in a given region.
    \item \textbf{Geospatial Grid Analysis}\\
      Compare regions to find the most sustainable locations
      for pretraining, given their grid mix.
  \end{itemize}

  \vspace{0.3cm}
  \textbf{Optional:}
  \begin{itemize}
    \item \textbf{Computational Granularity \& Resumption Logic}\\
      Study how checkpoint frequency (batch vs.\ epoch, etc.)
      impacts the trade-off between CO$_2$ savings and overhead.
  \end{itemize}

\end{frame}
